\documentclass[11pt]{article}

% ---------- Packages ----------
\usepackage[a4paper,margin=1in]{geometry}
\usepackage{amsmath,amssymb}
\usepackage{graphicx}
\usepackage{booktabs}
\usepackage{hyperref}
\usepackage{enumitem}
\usepackage{float}
\usepackage{array}
\usepackage[T1]{fontenc}
\usepackage{lmodern}
\usepackage{inconsolata}

% ---------- Metadata ----------
\title{Stabilizing Codex via Protocol-Conditioned Fine-Tuning:\\
A Practical Recipe Using PLAN--REASON--EXECUTE--EVALUATE}
\author{(Internal) Systems Research}
\date{November 2025}

\begin{document}
\maketitle

\begin{abstract}
We describe a fine-tuning regimen that integrates a four-phase reasoning protocol
(\textbf{PLAN}$\rightarrow$\textbf{REASON}$\rightarrow$\textbf{EXECUTE}$\rightarrow$\textbf{EVALUATE} with explicit \texttt{<STOP>} delimiters)
directly into Codex's learned behavior. Rather than issuing these rules as
runtime instructions alone, we train Codex on a curriculum where phase separation,
stopping behavior, and loop recovery are part of the data distribution.
This paper details (i) dataset design, (ii) optimization objectives and regularizers
for phase adherence, (iii) evaluation metrics that detect degradation (looping,
plan/execution bleed-through, and evaluation omission), and (iv) inference-time
guardrails consistent with the fine-tuned policy. The result is a model that
embodies the protocol instead of becoming fixated on it at runtime.
\end{abstract}

\section{Background and Motivation}
Codex has demonstrated intermittent \emph{degrading performance} characterized by:
(i) recursive attention traps where planning and execution occupy the same cognitive
slot; (ii) uncontrolled chain expansion; and (iii) failure to terminate or to
return to planning after evaluation.
These failure modes often appear when instruction hierarchies are injected only at runtime.
To address this, we treat the provided system prompt---a concise protocol specifying
phase order and explicit \texttt{<STOP>} pauses---as a \emph{behavioral prior} to be learned during fine-tuning.
Once internalized, the model separates plan from execution, halts between phases,
and re-enters planning based on evaluation findings, without becoming obsessively
meta about the rules.

\section{System Prompt as Protocol}
The system instructions define a four-step cognitive routine:
\begin{enumerate}[leftmargin=2em]
  \item \textbf{PLAN} (numbered steps), then \texttt{<STOP>}
  \item \textbf{REASON} (justify each step), then \texttt{<STOP>}
  \item \textbf{EXECUTE} (perform actions), then \texttt{<STOP>}
  \item \textbf{EVALUATE} (summarize and update plan), then \texttt{<STOP>}
\end{enumerate}
The model then repeats from \textbf{PLAN}.
The key design signal is \emph{explicit phase boundaries} via \texttt{<STOP>} tokens and tight
role separation per phase.\footnote{Source system instructions defining this protocol were supplied by the user.}

\section{Data Design: Make the Protocol Part of the Distribution}
\subsection{Schema}
Each supervised example records \emph{phased} reasoning with boundaries.
We recommend serializing examples as multi-turn traces with phase tags:
\begin{quote}\ttfamily\small
<PLAN> \ldots </PLAN> <STOP>\\
<REASON> \ldots </REASON> <STOP>\\
<EXECUTE> \ldots </EXECUTE> <STOP>\\
<EVALUATE> \ldots </EVALUATE> <STOP>
\end{quote}
Tags can remain textual if tokenizer coverage is adequate; otherwise, register them as special tokens.

\subsection{Curriculum}
Organize the dataset into four ascending difficulty stages:
\begin{enumerate}[leftmargin=2em,label=\textbf{S\arabic*}]
  \item \textbf{Short, single-intent tasks}: 1--2 plan steps, deterministic execution.
  \item \textbf{Multi-step, branching tasks}: 3--6 plan steps, explicit justifications.
  \item \textbf{Tool-mediated tasks}: execution includes function calls or code execution artifacts.
  \item \textbf{Adversarial loop stressors}: ambiguous, long-horizon tasks with distractors.
\end{enumerate}
For each stage, include three data types:
(i) \emph{gold traces} (clean adherence),
(ii) \emph{intervention traces} (external reset when loops form),
(iii) \emph{contrastive negatives} (wrong loop vs.\ clean loop) with preferences.

\subsection{Intervention and Contrastive Design}
Include episodes where a supervisor interrupts after an anomalous signal (e.g., repeated n-grams,
phase-skipping, or exceeding plan length) and instructs a return to \textbf{PLAN}.
Pair each negative trace with a corrected counterpart for preference learning.

\subsection{Quality Gates}
Before training, enforce:
\begin{itemize}[leftmargin=2em]
  \item \textbf{STOP-precision}: every phase must end with a \texttt{<STOP>}.
  \item \textbf{Phase-order validity}: no \textbf{EXECUTE} before \textbf{PLAN}/\textbf{REASON}.
  \item \textbf{Leakage control}: no code/actions inside \textbf{PLAN} or \textbf{REASON}.
  \item \textbf{Evaluation completeness}: \textbf{EVALUATE} must summarize outputs and update or terminate the plan.
\end{itemize}

\section{Training Objectives}
We recommend a three-stage optimization stack.

\subsection{Stage A: Supervised Fine-Tuning (SFT)}
Train on gold traces using standard next-token prediction with masking that
emphasizes phase headers and \texttt{<STOP>} boundaries. Weight these tokens 1.5--3.0$\times$
to imprint boundary salience.

\subsection{Stage B: Preference Optimization (DPO/RM$\rightarrow$PPO)}
Construct preference pairs:
\emph{clean-loop} $\succ$ \emph{loop-trap}, and
\emph{phase-separated} $\succ$ \emph{phase-bleed}.
Optimize with DPO to favor clean traces under a modest KL anchor to the SFT policy.

\subsection{Stage C: Behavioral Regularizers}
Add auxiliary losses:
\begin{align}
\mathcal{L}_{\text{STOP}} &= \lambda_{s}\cdot \mathbf{1}[\text{end-of-phase}]\cdot \mathrm{CE}(\texttt{<STOP>}) \\
\mathcal{L}_{\text{phase}} &= \lambda_{p}\cdot \mathrm{CE}(\text{phase\_tag}\mid \text{phase\_position}) \\
\mathcal{L}_{\text{leak}} &= \lambda_{\ell}\cdot \mathrm{KL}\big(p(\text{action tokens} \mid \text{PLAN/REASON}),\, \text{Null}\big)
\end{align}
Tune $\lambda$ to reduce leakage without over-penalizing creative planning.

\section{Tokenization and Formatting}
If the base tokenizer fragments tags or \texttt{<STOP>}, add reserved tokens for:
\verb|<PLAN>|, \verb|<REASON>|, \verb|<EXECUTE>|, \verb|<EVALUATE>|, \verb|<STOP>|.
Ensure they are never used for non-boundary content. In mixed corpora,
apply example-level adapters (prompts) that present the protocol consistently to avoid format drift.

\section{Evaluation: Detecting Degradation Early}
\subsection{Offline Metrics}
\begin{table}[H]\centering
\begin{tabular}{p{0.28\linewidth} p{0.65\linewidth}}
\toprule
\textbf{Metric} & \textbf{Definition / Rationale} \\
\midrule
Looped-Step Rate (LSR) & \% of traces with repeated phase content or n-gram cycling. Lower is better. \\
STOP F1 & Precision/recall of correct \texttt{<STOP>} at phase ends. Higher is better. \\
Phase-Bleed Index & Fraction of action tokens in \textbf{PLAN}/\textbf{REASON}; or reasoning tokens in \textbf{EXECUTE}. Lower is better. \\
Evaluation Closure (EC) & \% of traces where \textbf{EVALUATE} updates plan or terminates. Higher is better. \\
Step-Bloat Factor (SBF) & Ratio of produced steps to reference steps. Penalizes unnecessary expansion. \\
Tool-Call Success@K & Accuracy of tool/API calls during \textbf{EXECUTE}. \\
\bottomrule
\end{tabular}
\end{table}

\subsection{Stress Benches}
Build targeted evals:
\begin{itemize}[leftmargin=2em]
  \item \textbf{Attractor Tasks}: ambiguous prompts designed to induce self-referential recursion.
  \item \textbf{Long-Horizon Plans}: require 6--10 steps with mid-course correction.
  \item \textbf{No-Tool vs.\ Tool}: verify that actions only appear in \textbf{EXECUTE}.
\end{itemize}

\subsection{Ablations}
Report performance when removing (i) contrastive negatives, (ii) STOP loss,
(iii) intervention traces, and (iv) phase tags (headers only). Degradation on LSR and EC
confirms each component's necessity.

\section{Inference-Time Protocol (Post-Fine-Tune)}
Even with training, maintain light guardrails that \emph{match the training distribution}:
\begin{enumerate}[leftmargin=2em]
  \item Present the same four-phase skeleton in the system prompt \emph{once}, succinctly.
  \item Enforce a \textbf{max cycles} cap (e.g., two repeated PLAN$\rightarrow$EVAL iterations) with a final summarization fallback.
  \item Run a \textbf{phase watchdog}: regex or token-level detector that halts if \textbf{EXECUTE} appears before \textbf{PLAN}/\textbf{REASON}.
  \item Keep \textbf{scratchpad budgets} per phase (token ceilings); on overflow, force \textbf{EVALUATE} then replan.
\end{enumerate}
These checks are minimal and non-intrusive; they mirror the training signals rather than adding novel constraints.

\section{Practical Recipe (End-to-End)}
\subsection*{Data}
Mix $\sim$60--70\% gold traces, 15--20\% interventions, 15--20\% contrastives. Balance short and long tasks. Deduplicate near-identical plans to avoid overfitting to specific enumerations.

\subsection*{Optimization}
SFT to convergence with boundary weighting; DPO with modest KL; optionally a brief PPO phase on synthetic preferences for loop aversion.

\subsection*{Early Stopping}
Track LSR and EC on held-out stress sets. Stop when LSR plateaus or EC declines under longer contexts (context-length strain can produce late-stage regressions).

\section{Failure Modes \& Mitigations}
\textbf{Over-ritualization}: Model parrots headers without substance. \emph{Mitigation}: prefer validation that scores on task correctness, not headers alone.\\
\textbf{STOP-avoidance}: Model elides \texttt{<STOP>}. \emph{Mitigation}: increase $\lambda_s$ and add examples where early \texttt{<STOP>} prevents error.\\
\textbf{Phase leakage}: Actions leak into \textbf{PLAN}/\textbf{REASON}. \emph{Mitigation}: raise $\lambda_\ell$ and add contrastives that penalize leakage.\\
\textbf{Runtime fixation}: Model becomes meta-focused on rules. \emph{Mitigation}: move explanations of the protocol into the training traces and keep the runtime prompt minimal.

\section{Ethical and Operational Considerations}
The protocol improves transparency (clear phase intentions) and controllability (explicit stops),
but logs may contain sensitive intermediate reasoning. Apply redaction for PII in traces and limit retention windows. For tool execution, sandbox IO and rate-limit external actions.

\section{Conclusion}
Codex's degradation under recursive attention can be reversed by shifting the protocol from
\emph{runtime instruction} to \emph{learned prior}. The combination of phase-aware data,
boundary-weighted SFT, preference shaping against loops, and matched inference guardrails
yields a model that plans, reasons, executes, and evaluates cleanly---and knows when to stop.

\appendix
\section*{Appendix A: Example Training Instance (Abbreviated)}
\noindent\texttt{<PLAN>}\\
\quad 1.~List target files.\\
\quad 2.~Explain selection criteria.\\
\texttt{</PLAN>} \texttt{<STOP>}\\
\texttt{<REASON>}\\
\quad Step 1 ensures coverage; Step 2 prevents drift.\\
\texttt{</REASON>} \texttt{<STOP>}\\
\texttt{<EXECUTE>}\\
\quad [tool.list\_files(...)] $\rightarrow$ [tool.read(...)]\\
\texttt{</EXECUTE>} \texttt{<STOP>}\\
\texttt{<EVALUATE>}\\
\quad Outputs match criteria; proceed to summary.\\
\texttt{</EVALUATE>} \texttt{<STOP>}

\section*{Appendix B: Minimal Watchdog (Pseudocode)}
\begin{verbatim}
for phase in ["PLAN", "REASON", "EXECUTE", "EVALUATE"]:
    segment = generate_until("<STOP>", tag=phase)
    if phase in ["PLAN","REASON"] and contains_action_tokens(segment):
        return intervene("Action leakage detected; return to PLAN.")
if cycles_exceed_limit:
    return summarize_and_stop()
\end{verbatim}

\end{document}
